% 05-rasprava.tex — Discussion (TASK §F.6, §D.4).
\section{Rasprava}
\label{sec:rasprava}

\subsection{Kašnjenje i osjećaj kontrole}

Niz mjerenja s dodanim kašnjenjem pokazuje jasan prijelaz u načinu vođenja koji ovisi o ukupnom vremenskom odzivu sustava. Prijelaz je zabilježen u ovom postavu, s jednim operaterom i jednim promatračem, pa opisuje ponašanje ovog sustava, a ne populacijski prag. Do 50 ms dodanog kašnjenja (ukupno fizikalno kašnjenje do $\approx 435$ ms) operater vodi robota kontinuirano; blago kašnjenje već primjećuje, ali sustav ocjenjuje potpuno upravljivim. Pri 100 ms dodanog kašnjenja (ukupno $\approx 485$ ms, odnosno oko pola sekunde) operater počinje svjesno usporavati i čekati odziv robota u točkama promjene smjera. Pri 200 ms i 400 ms dodanog kašnjenja, kada ukupno kašnjenje doseže 585 ms i čak 785 ms (gotovo cijelu sekundu zaostajanja), operater u potpunosti prelazi na diskontinuiranu strategiju pomaka i čekanja (\emph{move-and-wait}), pri kojoj izvede diskretan pomak ruke, pričeka vizualnu potvrdu robota, provjeri poziciju te tek potom planira sljedeći pokret. Promatrač je isti prijelaz opisao kao trenutak u kojem operater počinje ulagati znatan kognitivni napor u kompenzaciju tromosti sustava umjesto u samo obavljanje zadatka.

Taj je nalaz u skladu s literaturom o teleoperaciji. Louca i suradnici pri dvosmjernim kašnjenjima od 0 do 2600 ms bilježe sustavan porast udjela čekanja i pad uspješnosti kroz sve mjerene metrike, a prag od približno 500 ms iznad kojeg pogreške naglo rastu pripisuju ranijem mjerenju Pereza i suradnika \cite[str.~14]{louca2024}. Njihova su kašnjenja dvosmjerna, a ovdje mjerena jednosmjerna, pa se pragovi ne uspoređuju izravno. Budući da sustav već u baznom radu unosi 385 ms dinamičkog odziva (zbog One-Euro filtra, sigurnosnog sloja i regulacije pogona), dodavanje 100 do 200 ms pomiče ga u područje u kojem oba rada bilježe prijelaz na diskontinuirano vođenje, dok pri dodanih 400 ms (ukupno 785 ms) kontinuirano vođenje u stvarnom vremenu postaje nemoguće. Wang i suradnici pokazuju da pod promjenjivim kašnjenjem stabilnost i transparentnost ulaze u kompromis \cite[str.~5]{wang2025}. U našem radu regulacijski krug pri 785 ms ukupnog odziva zadržava potpunu numeričku stabilnost bez pojave oscilacija, no zbog prevelikog vremenskog pomaka operater gubi mogućnost izravnog intuitivnog vođenja.

Praktična posljedica za projektiranje pHRI sustava jest da praćenje računalne latencije nije dovoljno. Latencija proračuna upravljačke petlje iznosi svega 3 ms, dok ukupno kašnjenje koje operater doživljava iznosi 385 ms. Razlika je gotovo u cijelosti posljedica faznog pomaka filtra i sigurnosnog sloja. Komparativna simulacija u okruženju URSim potvrđuje ovaj uvid jer isti zapis reproduciran u simulatoru zaostaje 248 ms, dakle 137 ms manje nego na fizičkom robotu. Ta razlika mjeri odmak između dvaju okruženja, a ne komponentu budžeta stvarnog lanca, čije razlaganje daje tablica~\ref{tab:budzet}. Sustav koji deklarira samo računsku latenciju izgleda stotinjak puta brži nego što se stvarno ponaša u interakciji.

\subsection{Koje su pogreške prihvatljive}

Dvije vrste odstupanja mjerene u ovom radu operater doživljava potpuno
različito. Fazni pomak od nekoliko stotina milisekundi mjerljivo je velik, ali
dok je konstantan, čovjek ga nauči predviđati i kompenzirati unaprijed, što
se i potvrdilo pri manjim iznosima dodanog kašnjenja. Amplitudno odstupanje i
diskontinuitet bitno su nepovoljniji za korisnika. Rad bez filtra imao je objektivno najmanju pogrešku,
ali je subjektivno ocijenjen kao nesiguran zbog trzanja motora. Prigušenje koje One-Euro filtar osigurava
stvara gladak pokret koji ulijeva povjerenje operateru.

To opravdava izbor filtra One-Euro kao primarnog rješenja jer postiže najbolji omjer glatkoće i brzine odziva, uz najmanji broj zasićenja brzine među filtriranim varijantama. Martini i suradnici uz sirovu pogrešku izvještavaju i pogrešku akceleracije kao zasebnu mjeru podrhtavanja \cite[str.~5]{martini2024}, čime pokazuju da glatkoća zaslužuje vlastitu metriku. Ovdje izmjereni kompromis ide korak dalje jer dvije metrike pokazuju u suprotnom smjeru, pa izbor filtra postaje odluka o tome koja od njih vodi.

\subsection{Granica prema autonomiji}

Sustav radi kao teleoperacijski postav u kojem robot nema vlastiti model zadatka već izvršava naredbe izvedene iz položaja ruke. Stroppa i suradnici opisuju dijeljeno upravljanje kao spektar između ručnog vođenja i autonomije \cite[str.~5]{stroppa2023}. Naš sigurnosni sloj predstavlja prvi korak prema tom spektru jer robot autonomno korigira naredbu kada ona prelazi zadane fizikalne granice. Pri ispitivanju brzine ta je korekcija trajala gotovo četvrtinu pokusa, a operater ju je opisao kao ugodno i prirodno ponašanje stroja.

To se izravno povezuje s pitanjem osjećaja djelovanja (\emph{sense of agency}) u interakciji čovjeka i robota \cite[str.~1]{glawe2026}, pri čemu operater zadržava osjećaj kontrole dok su intervencije robota predvidljive i u skladu s fizikalnom intuicijom (poput glatkog ograničenja brzine). Nasuprot tome, nagli prekidi ili zamrzavanja pojedinih zglobova pri okluziji markera doživljavaju se kao smetnja u radu.

\subsection{Ograničenja}

Rezultate rada treba interpretirati uz sljedeća poznata ograničenja.

\textbf{Subjektivna procjena ima dva sudionika.} Operater i promatrač ispunili su strukturirani upitnik u slobodnoj formi. Znanstvena korisnička studija sa statističkim uzorkom i validiranim psihometrijskim ljestvicama zahtijeva veći broj ispitanika, što nadilazi opseg ovog rada. Opažanja operatera stoga služe kao kvalitativna potvrda uz izmjerene objektivne teleometrijske veličine.

\textbf{Markeri prate odjeću, a ne kostur.} Markeri su postavljeni na sportsku majicu i rukavicu, pa mjerenje uključuje i elastični pomak tkanine i mekog tkiva. Korištenje krutih ortoza uklonilo bi taj pomak, ali bi narušilo prirodnost ljudskog pokreta. Analiza osjetljivosti pokazuje da pomak markera od 1 mm unosi pogrešku kuta od 0,77 do 1,40$^\circ$, što je znatno manje od raspona pokreta zgloba.

\textbf{Efektivna frekvencija ulaznog toka podataka.} Uslijed mrežnog grupiranja paketa, ulazni signali stizali su u petlju frekvencijom od 31 do 73 Hz umjesto nazivnih 120 Hz. One-Euro filtar je stoga radio uz efektivno dulji period uzorkovanja, što objašnjava izmjereni fazni pomak od 200 ms.

\textbf{Mjerenje latencije obrade.} Vremenska oznaka postavlja se pri primitku paketa na upravljačko računalo, pa izmjerene 3 ms obuhvaćaju obradu i slanje, dok je optički put kroz Motive procijenjen iz stabilnosti okvira na oko 8,4 ms.

\textbf{Statičko ograničenje radnog prostora i brzine.} Granice brzine i kuta postavljene su fiksno, bez dinamičkog mjerenja udaljenosti čovjeka od robota. Fiksna granica vrijedi u cijelom radnom prostoru, pa je pri najduljem kraku najmanje konzervativna; PFL proračun u odjeljku~\ref{sec:postav} pokazuje da je zbog toga elevacija ramena jedina os koja pri punom dosegu prelazi prag norme ISO/TS 15066.

\FloatBarrier
